\section{Glossar und Abkürzungsverzeichnis (OKB)}

\begin{longtable}{p{0.25\textwidth} p{0.68\textwidth}}
\textbf{Begriff / Abkürzung} & \textbf{Erläuterung} \\
\hline
\endhead

OKB &
Bearbeitungskürzel des Projektmitglieds Buse Okcu. \\

PAA &
Bearbeitungskürzel des Projektmitglieds Arsenii Parhom. \\

TED &
Bearbeitungskürzel des Projektmitglieds Daniel Tempel. \\

Avisierungsdashboard &
Geschützter Bereich des Webfrontends für die Disposition. Dort werden Touren, Aufträge, Bestätigungsstatus, frühere Anfragen und Aktionen zur erneuten Avisierung angezeigt. \\

Avisierungsservice &
Backend-Komponente der digitalen Lieferavisierung. Der Service verarbeitet Avisierungsaufträge, steuert den Workflow, persistiert Zustände, versendet Benachrichtigungen und stellt Schnittstellen bereit. \\

Webfrontend &
React-basierte Benutzeroberfläche der digitalen Lieferavisierung. Sie umfasst das Avisierungsdashboard für Disponenten und die Bestätigungsseite für Kunden. \\

Bestätigungsseite &
Kundenseitige Oberfläche zur Prüfung der Lieferinformationen sowie zur Bestätigung oder Ablehnung des vorgeschlagenen Liefertermins. \\

E2E-Test &
End-to-End-Test; Test eines vollständigen Benutzer- oder Systemablaufs über mehrere Komponenten hinweg. \\

CSRF &
Cross-Site Request Forgery; Angriffsmuster, bei dem ungewollte zustandsändernde Anfragen ausgelöst werden. Dashboard-Anfragen werden deshalb mit einem CSRF-Token geschützt. \\

DISPO &
Externes führendes System für Auftrags-, Tour- und Lieferplanungsdaten. Die Liefer- und Tourenplanung verbleibt in DISPO. \\

Geocoding &
Umwandlung einer textuellen Adresse in geografische Koordinaten. Im Projekt wird dies für die Zielposition der Lieferverfolgung verwendet. \\

Nominatim &
Geocoding-Dienst auf Basis von OpenStreetMap-Daten, der im Projekt zur Ermittlung von Zielkoordinaten verwendet wird. \\

Twilio &
Externer Messaging-Dienst, der im Projekt für den Versand von SMS- und WhatsApp-Nachrichten eingesetzt wird. \\

\end{longtable}
